\section{Inference strategies}
\label{app:inference}

We performed various experiments to measure model quality for given runtime cost. 
The runtime can be tuned by changing the number of processed patches, or by using choosing different size of the model.

We firstly looked at how model quality changes in respect to decreasing area of the image compared to native resolution, presented in Figure~\ref{fig:inference2}. 
We observed that on ImageNet~\cite{deng2009imagenet} model retains most of the quality down to 40\% of the image size. 
After that, the quality drastically decreases.
On the other hand, increasing the size of the image have a diminishing return in quality. 
This can be directly compared with random token dropping as an alternative to resizing, which showed to be very ineffective way to decrease number of patches during inference - Figure~\ref{fig:inference_drop}.

Please note that this highly depends on the native resolution of the images in the dataset - e.g. dataset with twice as big images than ImageNet can probably be safely resized to 20\% of area. 


\begin{figure}[h]
\vspace{-5pt}
\centering
\begin{subfigure}[t]{0.33\textwidth}
\includegraphics[width=\textwidth]{figs/imagenet_area_resize.pdf}
% \caption{The effect of resizing the image}
\caption{}
\label{fig:inference2}
\end{subfigure}
% \hfill
\begin{subfigure}[t]{0.33\textwidth}
\includegraphics[width=\textwidth]{figs/imagenet_area_resize_drop.pdf}
\caption{}
% \caption{Dropping random tokens is ineffective way to decrease number of patches compared to resizing the image.}
\label{fig:inference_drop}
\end{subfigure}
% \hfill
\begin{subfigure}[t]{0.33\textwidth}
\includegraphics[width=\textwidth]{figs/imagenet_sequence_len_no_upscale.pdf}
\caption{}
% \caption{Given number of patches as compute budget, it is beneficial to upscale the image}
\label{fig:inference5}
\end{subfigure}
\caption{(a) The effect of resizing the image. (b) Dropping random tokens is ineffective way to decrease number of patches compared to resizing the image. Data from \navit-L/16. (c) Given number of patches as compute budget, it is beneficial to upscale the image.}
\vspace{-5pt}
\end{figure}

A better way to quantify the performance is by giving a constant compute budget corresponding to number of patches. Figure~\ref{fig:inference4} shows that resizing the image (while preserving aspect ratio) to 256 tokens retains most of the quality (within 0.3\%). 
This corresponds to 256x256 area (given patch size of 16). 
At 128 tokens (181x181 area) the quality difference reaches 1\% and drops significantly after that.

Here we also resized the image past its native resolution in case it already fit the given sequence length budget. 
We observed that it is beneficial to resize the image to the given sequence length past the native resolution to keep monotonic increase in quality, which is showed on
Figure~\ref{fig:inference5}.

Figure~\ref{fig:inference3} presents the runtime of \navit-L/16 and \navit-B/16 for different sequence lengths. We can see that \navit-L/16 at sequence length 128 is as fast as \navit-B/16 with sequence length 512, while having almost 1\% difference in quality.

\begin{figure}[t]
\vspace{-5pt}
\centering
\begin{subfigure}[t]{0.49\textwidth}
\includegraphics[width=\textwidth]{figs/imagenet_sequence_len.pdf}
% \caption{Quality on ImageNet in respect to number of patches (sequence length)}
\caption{}
\label{fig:inference4}
\end{subfigure}
\begin{subfigure}[t]{0.49\textwidth}
\includegraphics[width=\textwidth]{figs/imagenet_runtime_accuracy.pdf}
\caption{}
% \caption{Runtime of models compared to the accuracy on ImageNet}
\label{fig:inference3}
\end{subfigure}
\caption{(a) Quality on ImageNet in respect to number of patches (sequence length). (b) Runtime of models compared to the accuracy on ImageNet.}
\vspace{-5pt}
\end{figure}


\section{Cascades}
\vspace{-5pt}
Another strategy to be more compute efficient would be to assign more tokens (and thus FLOPs) to the examples deemed hard by the model.
This is in particular interesting for bulk inference workloads, where one can amortize over large datasets and where only the total inference time matters.
\begin{figure}[h]
\vspace{-5pt}
\centering
\includegraphics[width=.5\textwidth]{figs/cascades.pdf}
\caption{Performance of a model cascade versus the average inference time. The labels at the select points denote the number of tokens at that scale.}\label{fig:cascades}
\vspace{-5pt}
\end{figure}

To evaluate the feasibility of this approach, we consider two sequence lengths $n_1$ and $n_2$ with respective inference times $t_1$ and $t_2$.
Then, we (i) send all examples though the model with $n_1$ tokens, and send only the $\alpha\in(0,1)$-fraction deemed hardest (those with the smallest maximum probability) to the model with $n_2$ tokens.
To have an input almost exactly fit into $n_1$ or $n_2$ tokens we perform and aspect ratio preserving resize.
Hence, the total amortized inference time per-example is $t_i + \alpha t_2$, while the accuracy obtained by combining the accuracy of the first model on the $1-\alpha$ most-confident fraction of the data, and the performance of the more expensive model on the remaining data.
By considering several pairs of models and varying $\alpha$ we obtain the plot in \Cref{fig:cascades}.
As we can see this strategy is indeed useful and provides not only the best performing models at given compute budgets, but because $\alpha$ is a real parameter one can obtain very fine-grained trade-offs.

\section{Calibration}
\label{app:calibration}

To evaluate behaviour of the predicted uncertainties with scale, we compute the calibration error of a -B sized ImageNet-finetuned model (the -L model performs similarly).
Note that these models were trained with sigmoid loss, i.e., the 1000 labels were predicted independently without enforcing that the probabilities should sum up to 1.
As we varied the sequence of tokens per example between 128 and 1024, we obtained very stable calibration errors (top-1, using $\ell_1$ and $30$ buckets, i.e., the settings from \cite{nixon2019measuring}), which we present in \Cref{tbl:calibration}.

\begin{table}[h]
\centering
\caption{Expected calibration error on ImageNet-1K with varying sequence lengths.}
\label{tbl:calibration}
\begin{tabular}{r | c c c c c c c}
\toprule
Sequence Length & 128 & 256 & 384 & 512 & 640 & 768 & 1024 \\
\midrule
Calibration Error & 0.047 & 0.046 & 0.048 & 0.047 & 0.047 & 0.046 & 0.045 \\
\bottomrule
\end{tabular}

\end{table}


